\Askhsh{27318}{2}
\begin{ylh}{1}
    \item 1.2 Συναρτήσεις
    \item 1.4 Όριο συνάρτησης στο Χο
    \item 1.6 Μη πεπερασμένο όριο στο Χο
    \item 1.8 Συνέχεια συνάρτησης.
\end{ylh}
\centering
\begin{tikzpicture}
\begin{axis}[
    axis lines=middle,
    xmin=-0.5, xmax=14.5,
    ymin=-1.5, ymax=6.5,
    xtick={0,1,2,3,4,5,6,7,8,9,10,11,12,13,14},
    ytick={-1,0,1,2,3,4,5,6},
    grid=both,
    grid style={gray!30},
    tick label style={font=\scriptsize},
    width=0.8\textwidth, height=0.55\textheight, % Adjust size for better fit and readability
]
    % Dashed lines for specific y-values
    \addplot[dashed, secondarycolor] coordinates {(-0.5,2) (14.5,2)};
    \addplot[dashed, secondarycolor] coordinates {(-0.5,4) (14.5,4)};
    \addplot[dashed, secondarycolor] coordinates {(-0.5,6) (14.5,6)};

    % Dashed lines for specific x-values
    \addplot[dashed, secondarycolor] coordinates {(4,-1.5) (4,6.5)};
    \addplot[dashed, secondarycolor] coordinates {(8,-1.5) (8,6.5)}; % This dashed line is not in the original image but implied by the peak
    \addplot[dashed, secondarycolor] coordinates {(9,-1.5) (9,6.5)};

    % Plotting the function segments
    % Segment 1: from (0,-1) to (4,0) passing through a local maximum
    \addplot[very thick, maincolor, smooth] coordinates {(0,-1) (1,0) (3,1) (4,0)};
    
    % Segment 2: from (4,2) (open circle) to (6,0)
    % We approximate the curve starting just after x=4, coming from y=2, and going to (6,0).
    \addplot[very thick, maincolor, smooth] coordinates {(4.001, 1.99) (5,1) (6,0)}; 
    
    % Segment 3: from (6,0) to (9,4) (open circle) passing through a local maximum
    % We approximate the curve ending just before x=9, approaching y=4.
    \addplot[very thick, maincolor, smooth] coordinates {(6,0) (7,3) (8,5.8) (8.999, 4.01)}; 
    
    % Segment 4: from (9,4) (open circle) to (14,4) (solid circle) passing through a local minimum
    % We approximate the curve starting just after x=9, coming from y=4.
    \addplot[very thick, maincolor, smooth] coordinates {(9.001, 3.99) (11,5) (13,2.7) (14,4)}; 

    % Solid circles (f(x) values)
    \addplot[only marks, mark=*, mark size=2.5pt, black] coordinates {(0,-1) (4,0) (6,0) (9,2) (14,4)};
    
    % Open circles (discontinuities/jump points)
    \addplot[only marks, mark=o, mark size=2.5pt, black] coordinates {(4,2) (9,4)};

\end{axis}
\end{tikzpicture}
\vspace{0.5cm}

Στο παρακάτω σχήμα δίνεται η γραφική παράσταση μιας συνάρτησης $f$. Γνωρίζουμε ότι η $f$ παίρνει θετικές τιμές κοντά στο έξι και ο οριζόντιος άξονας εφάπτεται στη γραφική της παράσταση στο σημείο αυτό.

\begin{alist}
\item Να βρείτε το πεδίο ορισμού και το σύνολο τιμών της. \monades{06}
\item Να εξετάσετε αν υπάρχουν και να βρείτε τα παρακάτω όρια:
    \begin{rlist}
    \item $\lim_{x \to 0} f(x)$
    \item $\lim_{x \to 4} f(x)$
    \item $\lim_{x \to 9} f(x)$
    \item $\lim_{x \to 14} f(x)$
    \item $\lim_{x \to 6} \dfrac{1}{f(x)}$
    \end{rlist}
    Για τα όρια που δεν υπάρχουν να αιτιολογήσετε την απάντησή σας. \monades{12}
\item Να βρείτε τα σημεία στα οποία η $f$ δεν είναι συνεχής. Να αιτιολογήσετε την απάντησή σας. \monades{07}
\end{alist}